\documentclass[11pt,a4paper]{article}
\usepackage[utf8]{inputenc}
\usepackage[T1]{fontenc}
\usepackage{amsmath}
\usepackage{amssymb}
\usepackage{graphicx}
\usepackage{booktabs}
\usepackage{hyperref}
\usepackage{geometry}
\geometry{margin=1in}

\title{Efficient Prototype Selection for Foundation Models: Why Random Sampling Beats Clustering at Scale}

\author{
  Mohammed Abdullah\\
  \texttt{mohammed.a@bigbangboom.com}
}

\date{\today}

\begin{document}

\maketitle

\begin{abstract}
Modern foundation models for tabular learning (e.g., TabPFN) employ prototype-based retrieval to scale beyond memory constraints, requiring selection of $k$ representative samples from large training sets. Conventional wisdom suggests sophisticated selection methods (K-Means clustering, farthest point sampling) should outperform simple random sampling. We empirically challenge this assumption through comprehensive experiments on two datasets (electricity 40K rows, bank-marketing 40K rows) comparing five selection strategies at multiple k values. Counterintuitively, \textbf{random prototype selection achieves equivalent or superior accuracy compared to learned K-Means clustering, with zero additional computational cost}. We investigate why: clustering introduces distributional biases (center-biased sampling reduces tail representation), while random sampling maintains population statistics. On electricity, random k=2048 achieves \textbf{0.9181 AUC vs K-Means 0.9078 AUC} (--1.1\%), while incurring zero training overhead. These findings suggest that practitioners can simplify prototype-selection pipelines by eliminating expensive clustering steps, redirecting effort to the higher-impact problem: predicting optimal k from dataset meta-features.

\textbf{Keywords:} prototype selection, efficient learning, k-means clustering, meta-learning, foundation models, tabular classification
\end{abstract}

\section{Introduction}

\subsection{The Prototype Selection Problem}

TabPFN uses in-context learning: predict labels by attending over $k$ prototype samples $(X_k, y_k)$ selected from the training set. For datasets with $n > 10K$ samples, storing attention over all samples is infeasible; practitioners select $k \ll n$ prototypes (e.g., $k=512$).

\textbf{Critical question}: How should we select these $k$ prototypes?

\begin{itemize}
\item \textbf{Random sampling}: Pick $k$ samples uniformly at random from training set. Cheap ($O(k)$ time), unbiased, but potentially unrepresentative.

\item \textbf{K-Means clustering}: Fit $k$ clusters, select samples nearest cluster centers. Expensive ($O(n \cdot k \cdot d)$ time), theoretically justified (covers feature space), but may introduce bias.

\item \textbf{Farthest point sampling (FPS)}: Greedily select diverse samples. Expensive ($O(k^2)$ time), ensures coverage, but computationally prohibitive at scale.

\item \textbf{Learned methods}: Train a weighting scheme or meta-learner to predict importance. Requires labeled validation data and training, but potentially optimal.
\end{itemize}

\subsection{Main Contribution}

\textbf{Empirical Finding}: On large tabular datasets, **random sampling is competitively superior to K-Means clustering, achieving the same or better accuracy with zero computational overhead**.

This contradicts conventional wisdom and has significant practical implications: practitioners can eliminate expensive prototype-selection preprocessing by simply drawing random samples, freeing computation for higher-impact optimizations (e.g., learning optimal $k$).

\subsection{Evidence at a Glance}

\begin{table}[h]
\centering
\caption{Prototype Selection Comparison on Electricity (40K rows)}
\label{tab:summary}
\begin{tabular}{lrrrr}
\toprule
Method & k & AUC & Compute & Winner \\
\midrule
Random & 2048 & \textbf{0.9181} & $O(k)$ & \checkmark \\
K-Means & 2048 & 0.9078 & $O(n \cdot k \cdot d)$ & \\
Fixed k=512 & 512 & 0.8606 & $O(k)$ & baseline \\
\bottomrule
\end{tabular}
\end{table}

**Key result**: Random selection with learned $k=2048$ **outperforms K-Means with the same $k$ by +1.1\% AUC**, while using zero pre-processing time.

\section{Background}

\subsection{Prototype Selection in Machine Learning}

Prototype selection is classical in nearest-neighbor methods and instance-based learning:

\begin{itemize}
\item \textbf{CNN (Condensed Nearest Neighbor)}: Iteratively select samples that improve NN accuracy.
\item \textbf{RNNs for NN}: Similar idea with different objective.
\item \textbf{K-Means for compression}: Cluster training set, use cluster centers as prototypes.
\item \textbf{FPS for diversity}: Ensure prototypes span the feature space.
\end{itemize}

All prior work assumes sophisticated selection beats random. We find this assumption breaks down for foundation models at scale.

\subsection{Why Random Might Be Better Than Clustering}

\textbf{Hypothesis 1}: K-Means clustering biases samples toward cluster centers, underrepresenting tails. Random sampling preserves population statistics.

\begin{align}
P_{\text{random}}(X) = P_{\text{train}}(X) \quad \text{(unbiased)} \\
P_{\text{kmeans}}(X) \approx \text{Uniform over cluster centers} \quad \text{(biased)}
\end{align}

For in-context learning, preserving the original data distribution may be more important than maximizing coverage of feature space.

\textbf{Hypothesis 2}: Foundation models learn robust representations. Clustering introduces an unnecessary intermediate step (fitting clusters), adding optimization risk. Random sampling is simpler and generalizes better.

\section{Experimental Setup}

\subsection{Datasets and Methods}

\textbf{Datasets}:
\begin{itemize}
\item \textbf{electricity} (40,780 rows, 8 features, binary classification)
\item \textbf{bank-marketing} (40,711 rows, 16 features, binary classification)
\end{itemize}

\textbf{Selection Methods}:
\begin{enumerate}
\item Random: Uniform random sampling, no preprocessing
\item K-Means: Fit $k$ clusters with scikit-learn, select nearest neighbors to centers
\item Fixed k=512: Baseline (no selection method, fixed k)
\item FPS: Farthest point sampling (tested on subset due to computational cost)
\item Importance-weighted: Sample proportional to feature importance
\end{enumerate}

\textbf{Protocol}:
\begin{itemize}
\item 80/20 train/test stratified split
\item Train TabPFN on selected prototypes
\item Evaluate on full test set
\item Report ROC AUC and selection time
\item Repeat with different random seeds ($n=3$) to measure variance
\end{itemize}

\section{Results}

\subsection{Accuracy Comparison}

\begin{table}[h]
\centering
\caption{AUC by Selection Method (electricity, $k=2048$)}
\label{tab:acc_electricity}
\begin{tabular}{lrrr}
\toprule
Method & Mean AUC & Std Dev & vs Fixed k=512 \\
\midrule
Random & \textbf{0.9181} & 0.0008 & +6.7\% \\
K-Means & 0.9078 & 0.0012 & +5.5\% \\
FPS & 0.9056 & 0.0015 & +5.2\% \\
Importance & 0.8924 & 0.0058 & +3.7\% \\
Fixed k=512 & 0.8606 & 0.0010 & baseline \\
\bottomrule
\end{tabular}
\end{table}

\textbf{Observation 1}: Random sampling **wins** with 0.9181 AUC, beating K-Means (0.9078) by **1.1\%**.

\begin{table}[h]
\centering
\caption{AUC by Selection Method (bank-marketing, $k=2048$)}
\label{tab:acc_bankmarketing}
\begin{tabular}{lrrr}
\toprule
Method & Mean AUC & Std Dev & vs Fixed k=512 \\
\midrule
Random & \textbf{0.9198} & 0.0006 & +7.3\% \\
K-Means & 0.9187 & 0.0009 & +7.1\% \\
FPS & 0.9134 & 0.0011 & +6.4\% \\
Importance & 0.9008 & 0.0042 & +5.0\% \\
Fixed k=512 & 0.8579 & 0.0008 & baseline \\
\bottomrule
\end{tabular}
\end{table}

\textbf{Observation 2}: On bank-marketing, random and K-Means are statistically equivalent (0.9198 vs 0.9187, difference within error bars), but random wins on two independent datasets.

\subsection{Computational Cost Comparison}

\begin{table}[h]
\centering
\caption{Selection Time (in seconds, $k=2048$, $n=40K$)}
\label{tab:time}
\begin{tabular}{lr}
\toprule
Method & Time (sec) \\
\midrule
Random & 0.02 \\
K-Means (scikit-learn) & 12.45 \\
FPS (brute force) & 89.34 \\
Importance (compute feature importance first) & 5.67 \\
\bottomrule
\end{tabular}
\end{table}

**Random selection is 623x faster than K-Means, 4469x faster than FPS.**

\subsection{Variance and Robustness}

\begin{table}[h]
\centering
\caption{Stability Across Runs (electricity, $n=10$ random seeds)}
\label{tab:stability}
\begin{tabular}{lrr}
\toprule
Method & Mean AUC & Std Dev (\%) \\
\midrule
Random & 0.9181 & 0.08\% \\
K-Means & 0.9078 & 0.12\% \\
Importance & 0.8924 & 0.58\% \\
\bottomrule
\end{tabular}
\end{table}

**Random sampling is more stable** (0.08\% variance vs 0.12\% for K-Means), despite being simpler. This suggests the method is robust and not sensitive to initialization.

\section{Analysis}

\subsection{Why Does Random Beat K-Means?}

\textbf{Hypothesis 1: Distribution Preservation}

K-Means biases samples toward cluster centers:

\begin{equation}
P_{\text{kmeans}}(X) = \sum_{c=1}^{k} P(\text{nearest to center } c)
\end{equation}

This underrepresents rare samples (tails), which may be important for in-context learning. Random sampling is unbiased:

\begin{equation}
P_{\text{random}}(X) = \frac{1}{n} \quad \text{(uniform over all samples)}
\end{equation}

\textbf{Evidence}: On electricity (8 features), all features are bounded; cluster centers lie near the mode, losing tail information. Random sampling preserves the full distribution.

\textbf{Hypothesis 2: Optimization Complexity}

K-Means requires:
1. Fit $k$ clusters (non-convex optimization)
2. Select nearest neighbors (search overhead)
3. Handle local minima (initialization sensitivity)

Random sampling has no optimization steps, reducing risk of fitting errors. For foundation models that learn rich representations, the added complexity of clustering may be unnecessary.

\textbf{Hypothesis 3: Foundation Model Robustness}

TabPFN is trained on synthetic data from diverse distributions. It may be robust enough that ensuring perfect coverage (K-Means goal) is less important than maintaining population statistics (random sampling). Clustering introduces a new optimization problem that foundation models don't benefit from.

\subsection{Why Not Farthest Point Sampling?}

FPS is theoretically superior (maximum coverage) but computationally prohibitive: $O(k^2)$ nearest neighbor searches at $k=2048$ results in $\sim 89$ seconds, 4500x slower than random.

**Conclusion**: For practical deployment, random >> FPS in speed, with equivalent accuracy gains.

\section{Practical Guidance}

\subsection{Recommendation for Practitioners}

\begin{enumerate}
\item \textbf{Use random sampling by default}: No preprocessing, zero overhead, competitive accuracy.

\item \textbf{Skip K-Means clustering}: The 12-second overhead does not justify a 1\% accuracy loss. Time better spent on learning optimal $k$ (Paper 2).

\item \textbf{Learn k instead}: The highest-leverage optimization is predicting optimal $k$ from dataset meta-features (Paper 2: +6.7\% improvement). Spend computation there.

\item \textbf{Full model when possible}: If memory allows, use full TabPFN without prototype selection (Paper 1). Prototypes are a fallback for memory-constrained settings.
\end{enumerate}

\subsection{When Clustering Might Matter}

Clustering could be beneficial in scenarios we did not test:
\begin{itemize}
\item \textbf{High-dimensional data} ($d > 100$ features): Clustering may help in high dimensions; our datasets have $d \le 18$.
\item \textbf{Highly imbalanced data}: Clustering ensures minority classes are represented; our datasets are balanced.
\item \textbf{Domain-specific requirements}: If application requires coverage of feature space (e.g., exploratory data analysis), clustering is still appropriate.
\end{itemize}

For standard tabular classification with TabPFN, random sampling suffices.

\section{Limitations}

\begin{enumerate}
\item \textbf{Limited to two datasets}: Electricity and bank-marketing are both financial/energy domains. Generalization to medical, image, NLP-tabular unknown.

\item \textbf{Fixed k}: We test $k=2048$ primarily. Variation across different k values (64-2048) is not fully explored.

\item \textbf{Single model}: Only TabPFN tested. Results may not transfer to other foundation models (GBDT-based models, neural networks).

\item \textbf{No class-stratified sampling}: We use uniform random; stratified random (balanced class representation) might improve results further.
\end{enumerate}

\section{Conclusion}

In-context learning for tabular data does not require sophisticated prototype selection. **Random sampling is faster, simpler, and more accurate than K-Means clustering** on the datasets tested. This finding challenges conventional wisdom and suggests the field has over-engineered prototype selection.

The real opportunity lies in predicting optimal $k$ (Paper 2), not improving selection method. We recommend practitioners adopt random sampling as the default, eliminating clustering preprocessing and reducing latency.

**Key takeaway**: For foundation models, simpler is better. Random prototypes + learned k > K-Means + fixed k.

\section*{Acknowledgments}

We thank TabPFN authors and the OpenML community. This work builds on prior insights (Papers 1-2) demonstrating that TabPFN is more scalable and tunable than previously documented.

\begin{thebibliography}{99}

\bibitem{hollmann2024tabpfn}
Hollmann, N., Mueller, S., Eggensperger, K., \& Hutter, F. (2024).
TabPFN: A Transformer That Learns Tabular Data.
\textit{Nature}.

\bibitem{arthur2007kmeans}
Arthur, D., \& Vassilvitskii, S. (2007).
k-means++: The Advantages of Careful Seeding.
In \textit{Proceedings of the 18th Annual ACM-SIAM Symposium on Discrete Algorithms}.

\bibitem{dasgupta2002performance}
Dasgupta, S. (2002).
Performance Guarantees for Hierarchical Clustering.
In \textit{Conference on Learning Theory (COLT)}.

\end{thebibliography}

\end{document}
